\documentclass{ltjsbook}
\usepackage{docmute} 
\usepackage{myStyle} 

\begin{document}
\VerbatimFootnotes

\chapter{よくある質問とミスの例}
\section{Frequently Miss and Comments}

Rmdでレポートを提出したのに，なんだか中身の問題じゃないのに突き返された，中身を見てくれよ！と思う人もいるかもしれません。テキストではRやRmdでの課題を提出するよう求めているところがありますが，その際よく見られる学生さんのミスとその対応についてのコメントをここにまとめておきます。

\subsection{FMC1；そもそもファイルの書式が違う}
Rmdで提出してください，Rで提出してください，という指示に対して，違うものが提出されてくることがあります。書式があっていない，というのは些細なことのように思えるかもしれませんが，学術論文は書式に捉われず内容に集中するためにも，書式は整えられたものである必要があります。学会誌に掲載されている論文も，みなさんが書く卒業論文も，レポートに至るまで，書式や指示に沿ったものを準備する必要があります。書式があっていない場合は，門前払いになっても文句が言えないのがアカデミックの世界です。

ということですので，RやRmdで提出してください，という指示があれば，RやRmdで提出してください。ファイルの種類は拡張子で分類され\footnote{拡張子についてはセクション\ref{suffix},Pp.\pageref{suffix}参照}，Rファイルは\verb|.R|，Rmdファイルは\verb|.Rmd|という拡張子になっています。たまに\verb|.Rproj|というファイルを提出してくる人がいますが，これはRプロジェクトのファイルで，これにはRスクリプトも文書も含まれておらず，みなさんの計算環境情報が少し書いてあるだけです。また，\verb|.Rmd|だけ入っていれば良いかというとそうではありません。まれに\verb|Filename.Rmd.R|というようなファイルを送ってくる人がいます。これはファイル名にピリオドが含まれているだけで，ファイルの種類を識別する拡張子は\verb|.R|ですからRmdファイルではありません\footnote{最近のOSは拡張子を表示しない設定になっているものも少なくないので，このようなミスが生じます。}。そもそもファイル名には2バイト文字はもちろん，\verb|!|や\verb|?|などの記号を含めるべきではありません。ピリオドも当然記号の一種ですから，ファイル名にするのは不適切です\footnote{長いファイル名などの場合，空白を入れるのも適切ではありません。できなくはないのですが，やるべきではないのです。どうしても空白を入れたければ，アンダースコアなどで区切ると良いでしょう。}。

ではどうやって適切なファイル形式にするか，ということですが，最も素直な方法としてはRStudioで新しくファイルを開くときに，R markdown形式(略してRmd)を選ぶようにしましょう。もし間違って，R script形式(\verb|.R|形式)で開いてしまった，というときは，そのファイルを破棄して新しく作り直すのが一番ですが，エディタペインの右下にあるファイル種類表示(図\ref{FileEditorPane})をクリックして修正することもできます。

\begin{figure}[ht]
    \centering
    \includegraphics[keepaspectratio,width=0.8\linewidth]{../../common_contents/figures/RstudioFileEditor.png}
    \caption{ファイルの種類を判別する方法}
    \label{FileEditorPane}
\end{figure}

\subsection{FMC2；やっぱりファイルの形式が違う}
Rmd形式のファイルは，拡張子だけで決まるものではありません。図\ref{RmdFileForm}にあるように，Rmdファイルは冒頭の6行(正確には\verb|---|で囲まれた領域)がYMALと呼ばれるところで，文書全体の設定をしています。その下に，Rのコードを実行する部分(\keyterm{チャンク}{ちゃんく}{chunk})や文章の領域があります。文章のところは\verb|#|記号で見出しを作ったりできます。

このYMAL部分が壊れている，あるいはチャンクが正しく記述されていない場合，拡張子が\verb|.Rmd|であっても適切なRmdファイルにはなっていません。
YMAL部分の書き方はよくわからない，という人も多いと思いますので，RStudioでRmdファイルを作ったときの状態をなるべく変更しないように注意すると良いでしょう\footnote{Rmdファイルを新しく開くときに，文書タイトルや著者名，日付，出力ファイル形式などを設定することができるウィンドウが開きますので，そこに必要な情報を書くと自動的にそれを使ったYMALが生成されます。また，本文として幾つかのサンプルコードが最初から含まれていますが，これらに関してはサンプルをそのまま使うことはないので，全て削除してしまって結構です。}。

\begin{figure}[ht]
    \centering
    \includegraphics[keepaspectratio,width=0.8\linewidth]{../../common_contents/figures/RmdFile.png}
    \caption{Rmdファイルの中身における書式}
    \label{RmdFileForm}
\end{figure}

チャンクはRのコードを書くところで，バッククォーテーション3つでくくるのが決まりです。チャンク領域が始まるところに\verb|{r}|とかいて「ここがRで計算するところですよ」というのを指定するわけです\footnote{ほかにもPythonやJuliaなど他の計算言語を混ぜることもできます。}。このバッククォーテーション3つ(\verb|```|)が全角だったり(\verb|｀|)，ダブルクォーテーションだったり(\verb|"|)すると，機械は正しくチャンクであると認識しません。フォントなどの見せかけ上は，微妙な違いのように見えますが，機械にとって違う文字列は違う意味を持ちますので注意してください\footnote{記号の名称や入力については，セクション\ref{signs},Pp.\pageref{signs}を参照してください。}。RStudioで編集する場合，チャンクの領域はやや灰色がかった強調表示がされますので，どこにチャンクがあるかわかると思います。もし強調表示されていないようであれば，チャンクとして認識されていない可能性を疑った方が良いでしょう。

チャンクはバッククォーテーション3つで開き，おわったら同じくバッククォーテーション3つで領域を閉じます。閉じなければずっとRの計算領域が続くと解釈されますし，最後までチャンクが閉じられないとRmdファイルとして正しくない書式ということになります。
チャンクが正しく入力できているかについて，常に注意を払っておく必要があります。また，自分でバッククォーテーション3つを使ってチャンク領域を開いたり閉じたりするのが面倒だ，という人はRStudioのチャンク挿入ボタンから挿入すると間違いないでしょう。

\begin{figure}[ht]
    \centering
    \includegraphics[keepaspectratio,width=0.8\linewidth]{../../common_contents/figures/InsertChunk.png}
    \caption{RStudioのチャンク挿入ボタン}
    \label{InsertChunk}
\end{figure}

\subsection{FMC3；Rmdがknitできない}
Rmdは文書作成に加えて，Rでの計算がセットになったファイル形式であり，文中の数字や分析結果は「書き写す」のではなく「その場で生成する」ものです。
生成する，というのはRmdファイルを変換して，HTMLやDOCX，PDF形式のファイルにすることを指します\footnote{HTMLはHyper Text Markup Languageの略で，ブラウザで開くファイル形式です。DOCXはMicrosoft Wordのファイル形式です。PDFはPortable Document Formatの略で，データを紙に印刷した状態のようにサイズを固定して出力したファイル形式であり，OSがもつビューワーやAdobe Acrobat Readerなどで見ることができます。}。このファイル変換を\keyterm{ニット}{knit}{knit}といいます。編み物を編むようなイメージですね。このときに，タイトルをつけ，見出しのサイズを変え，Rの計算をして結果を埋め込む作業をするわけです(図\ref{Knit})。こうすることで，結果のコピペを避けること，再現性を担保することができるようになるわけです。
\begin{figure}[ht]
    \centering
    \includegraphics[keepaspectratio,width=0.8\linewidth]{../../common_contents/figures/RmdKnit.png}
    \caption{knitしてレポートが完成する}
    \label{Knit}
\end{figure}
すでに説明したように，チャンクを使って計算に必要な指示をRmdファイルに書きます。この指示はエラーのないコードである必要があります。当然ですが，スペルミスや間違ったRコードは「実行できない」というエラーになるわけです。その場合，knitは中断され結果のファイルが出力されません。提出されたレポートはknitして出来上がったもののことを指しますから，knitできないRmdファイルを提出されてもレポートが提出されたことにはならないわけです。

knitできないRmdファイルになっていないか，というのはご自身のRStudio環境でknitして確認できることです。提出前に一度，きちんと機能するRmdファイルになっているかどうか確認するようにしてください。
以下で述べるように，自分の環境でknitできても，提出先の(私の)環境でknitできないということもあり得ます。しかし，自分の環境でknitできないのに提出先でできる，ということはあり得ませんので，「knitできませんよ」というコメントをつけて返される前に自分で確認するようにしてください。

\subsection{FMC4；外部環境を参照してしまう}

さて，Rで行う作業の中には，csvファイルを読み込むといった「外部のファイルを使う」指示もありえます。例えば次のようなコードです。
\begin{lstlisting}[language=c++,label={code:FAQ01},caption={Rmd中のRコードが外部環境を参照してしまう}]
dat <- read.csv("social_effects.csv")
dat <- read.table("clipboard")
\end{lstlisting}

このコードでは，上の行は\verb|social_effects.csv|というファイルを，下の行はクリップボードの内容を読み込むようになっています。ファイルを読み込む指示は，ファイルがなければ当然エラーになりますから注意が必要です。レポート等でRmdファイルを提出するとき，Rmdファイルの他に必要な読み込むべきファイルがあれば一緒に提出するようにしてください。
また，クリップボードの内容は再現できません。クリップボードとは，コピーアンドペーストのコピーを行ったとき，PCの内部で一時的にコピー内容を覚えておく場所のことです。つまり，みなさんがみなさんの環境で，クリップボード上にデータを一旦保持している場合は，このコードでエラーがしょうじることはありません。しかし，レポート提出先の(私の)環境で，事前にデータのコピー作業を行なっているわけではないのですから，提出先でエラーが発生します。

そもそもRmdファイルは，作業の再現性を担保するためのファイルになっているわけですから，クリップボードの利用のような「自分の環境だけで可能な記録されない作業」をそのファイルに含めるのは適切ではありません。
RmdファイルはRmdファイルだけで分析作業が完結するように，必要な記録は全て記載されている必要があります。同様に，分析作業に関係のない冗長な指示や無駄な指示はRmdファイルに含めるべきではありません。

\subsection{FMC5；提出先の環境を変更してしまう}

Rはさまざまなパッケージを使って分析環境を拡張することができます。パッケージは\keytermL{CRAN}{CRAN}を通じてインターネット経由で配布されますから，ネット環境があれば誰でも最新のパッケージをとってくることができます。
みなさんがRmdファイルの中で使うRの関数の中には，パッケージの関数も含まれているでしょう。パッケージがなければ関数が動きませんから，パッケージをインストールする作業もRmdに書いておきたいと思うかもしれません。しかし，これは推奨できません。

パッケージのインストールは，実行環境の準備にあたる作業です。Rmdファイルを使ってコードのやり取りをするとき，提出先の環境で分析することになりますが，提出先の環境にどのようなパッケージをいつ入れるかは，提出先の環境の管理人が判断すべき問題です。提出するRmdファイルにパッケージをインストールする関数，すなわち\verb|install.packages()|関数が含まれているというのは，相手の環境を勝手に操作してしまうことと同じであり，セキュリティ的にも適切な発想ではありません。

環境の準備は提出先の管理人が管理すべき問題であり，またすでにパッケージが入っている場合は無駄な上書き作業をさせることにもなります。また\verb|install.package|関数はCRANのサーバを参照したりしますから，適切な設定がなければRチャンク実行時にエラーが発生します。いずれにせよ，Rチャンクの中に\verb|install.package|関数を含めないようにしましょう。

以上がRmdファイルやRファイルでレポートを提出するときの留意点です。その他にもRやRStudioを使うときによくある質問がありますので，それらも一問一答型で紹介しておきましょう。


\section{Frequently Asked Questions；よくある質問と答え}
\label{FreaquentlyAskedQuestions}

\begin{FAQ}{ファイルを読み込もうとすると次のようなエラーが出ます。どうすれば良いですか。}{}
    \verb|'xxxxx'に不正なマルチバイト文字があります|
    \tcblower
    {\large A.} 文字化けの原因は，たとえばUTF-8形式で提供されているファイルを，Windows標準のCP932形式で読み込もうとする，という文字コードの不一致です。
    ですからそれをオプションで指定してあげれば問題解決です。
    \verb|read.csv|関数を使っている場合，オプションの指定は\verb|read.csv("foo.csv",fileEncoding="UTF-8")|のようにします。
    \verb|tidyverse|の\verb|read_csv|関数を使っている場合，localeオプションで，locale関数のencodingを明示的にUTF-8とします。\verb|read_csv("foo.csv", locale=locale(encoding = "UTF-8"))|のようにします。
\end{FAQ}

\begin{FAQ}{パッケージを読み込もうとすると次のようなエラーが出ます。どうすれば良いですか。}{}
    \verb|library(tidyverse)でエラー；'tidyverse'というパッケージはありません。|
    \tcblower
    {\large A.} パッケージがないので，インストールしてください。
\end{FAQ}


\begin{FAQ}{パッケージをインストールしようとすると次のようなエラーが出ます。どうすれば良いですか。}{}
    \begin{verbatim}
Warning in install.packages:
    'lib = "C:/Program Files/R/R-4.0.5/library"' is not writable
\end{verbatim}
    \tcblower
    {\large A.} ユーザがファイルに書き込みをする権限を持っていないので，(パッケージファイルをドライブに)書き込みできないというエラーです。
    RStudioを実行する時に，「管理者として実行」を選びましょう。
\end{FAQ}


\begin{FAQ}{パッケージをインストールしようとすると次のようなエラーが出ます。どうすれば良いですか。}{}
    \begin{verbatim}
Warning in install.packages:
    ディレクトリ 'C:〜(任意の文字列)〜\????' を作成できません。理由は```'Invalid argument'です。
\end{verbatim}
    \tcblower{\large A.}
    ユーザ名が全角文字を含んでいるため，文字化けしてR側から操作ができません。
    解決する方法は二つあります。
    \begin{description}
        \item [ユーザを作り直す方法] ユーザ名に全角文字を含まない，新しいユーザを作成します。設定＞アカウントから新しいアカウントを作りましょう。
        \item [インストールフォルダを指定する] Rのコンソールで\verb|.libPaths()|と実行すると，パッケージをインストールするフォルダが出てきますが，ここを変更する方法です。次の2ステップで対応できます。
              \begin{itemize}
                  \item まずCドライブのすぐ下にインストール先のフォルダを作ります。\verb|myLib|としましょう。
                  \item 次にRのコンソールで\verb|.libPaths("C:/myLib")|書いて実行します。
              \end{itemize}
              これでインストール先が変わりますので，書き込み・インストールができるようになります。老婆心ながら付け加えますと，フォルダ名はなんでも構いませんが，全角文字ではいけません。また場所もCドライブのすぐ下である必要はありませんが，全角文字や空白を含むフォルダの下に入れてしまってはいけません。OndDriveのようなクラウドを指定するのも良くありません(探しに行った時にオフラインだとまたエラーになります)。
    \end{description}
\end{FAQ}


\begin{FAQ}{ファイルが読み込めません}{}
    \begin{verbatim}
file(filename, "r", encoding = encoding) でエラー:
   コネクションを開くことができません
 追加情報:  警告メッセージ:
 file(filename, "r", encoding = encoding) で:
   ファイル 'foo.csv' を開くことができません: No such file or directory
    \end{verbatim}
    \tcblower{\large A.}
    指定された場所にファイルがないので，読み込むことができないエラーです。
    確認すべきは，「プロジェクトを開いているか」，「プロジェクトフォルダの中に当該ファイルはあるか」，「ファイル名のミススペルはないか」です。
    プロジェクトってなんだという人は，RStudioの基本に立ち戻り，プロジェクトでファイルやフォルダを管理するようにしてください。プロジェクト管理については，Pp.\pageref{UnderProject}にもその説明があります。

    プロジェクトによる管理とは要するに，Rが今見ているフォルダの場所を固定する方法です。プロジェクトを開くと，プロジェクトフォルダが「今見ているフォルダ(ワーキングディレクトリ)」になりますので，その中のファイルを参照することになります。
    プロジェクトフォルダの中に当該ファイルがないと読み込むことができませんので\footnote{フォルダの位置を相対的・絶対的に指定してやれば，どこにおいていても読み込むことはできますが，この問題で悩んでいる人は相対・絶対パスの指定というところでさらに疑問が深まることになると思いますので，気にしないでくれて結構です。相対・絶対パスが知りたい人は，付録\ref{computer},\pageref{setLocation}でファイル場所とは何かを再確認してください)。}，当該ファイルをプロジェクトフォルダ内に移してきてください。
\end{FAQ}

\begin{FAQ}{ANOVA君が読み込めません}{}
    \begin{verbatim}
file(filename, "r", encoding = encoding) でエラー:
   コネクションを開くことができません
 追加情報:  警告メッセージ:
 file(filename, "r", encoding = encoding) で:
   ファイル 'http:/riseki.php.xdomain.jp/index.php?plugin=attach&refer=
  ANOVA君&openfile=anovakun_486.txt' を開くことができません: Invalid argument
\end{verbatim}
    \tcblower{\large A.}
    ファイルを読み込みにいく先がURL，すなわちインターネット上になっています。ANOVA君のファイルを一度手元のPCのフォルダにダウンロードし，そのローカルのファイルの位置を指定して読み込むようにしてください。
\end{FAQ}

\begin{FAQ}{効果量として Hedges の g を算出してください。という指示について実行するとｇではなくｄが出てしまうようなのですが、どうすればよいでしょうか。コードは次の通りです。}{}
    \begin{verbatim}
cohen.d(value ~ condition, data = dat, hedges.crrection = T)
\end{verbatim}
    \tcblower{\large A.}
    Hedgesの補正(correction)のオプションが通っていません。スペルミスです。オプションのスペルが間違っているので無視されたので，関数名通りCohenのdが算出されます。
\end{FAQ}

% \begin{FAQ}{Rスクリプトをrmdファイルに変換する方法がわかりません。}{}
% \tcblower{\large A.}
% Rスクリプトを直接Rmdに変換する術はありません。Rmdファイルを作って、Rスクリプトはチャンクの中に書き込みます。Rmdやチャンクについてはセクション\ref{Rmarkdown},Pp.\pageref{Rmarkdown}を参照してください。
% \end{FAQ}

\begin{FAQ}{描画の際に次のような注意が出てきます。これはどういう意味ですか。}{}
    \begin{verbatim}Removed 671 rows containing missing values (geom_point).\end{verbatim}
    \tcblower{\large A.}
    「671件の行で欠損値が含まれています」ということです。つまりデータセットの中に欠損値(観測されていない，数値が入っていない行）が671件あったので，それは表示できませんでしたよ，という意味です。
    警告が嫌だということであれば，データセットの中で欠けているものを除外する必要があります。Rの関数では，\verb|na.omit|で欠損値を除外することができます。
\end{FAQ}

\begin{FAQ}{lm関数を実行するコードでオブジェクトがないと言われます。}{}
    \begin{verbatim}
result <- lm( weight - height, data = dat)
\end{verbatim}
    \tcblower{\large A.}
    従属変数と独立変数とをつなげるのはチルダ(\verb|~|)という記号です。そこがハイフン(\verb|-|)になっています。スペルミスの一種です。
\end{FAQ}

\begin{FAQ}{lm関数を実行するコードでオブジェクトがないと言われます。}{}
    \begin{verbatim}
result <- lm( weight - height )
\end{verbatim}
    \tcblower{\large A.}
    データを与えていないので，\verb|weight|というオブジェクトを探しに行って，見つからないというエラーです。\verb|data|オプションでデータを与えてあげてください。
\end{FAQ}

\end{document}